% ============================================
% DOC-27 (ZH)
% 中文翻译 — SPEC-935-R356
% ============================================
\section*{\textbf{DOC-27} --- 起源（1853）}
\markboth{DOC-27 --- 起源（1853）}{}

\noindent \textit{从鉴定档案中提取的文件。}

\vspace{0.5em}\noindent\rule{\textwidth}{0.5pt}\vspace{0.5em}


\noindent {\footnotesize\texttt{塞阿拉公共档案馆 $\cdot$ 散件文书 $\cdot$ 十九世纪}}\\[-0.2em]

\noindent\textbf{文书:} 若泽·伊纳西奥·德苏萨神父致塞阿拉省省长、公共救济委员会的信。\\
\noindent\textbf{日期:} 1853年3月14日。\\
\noindent\textbf{来源:} 雷梅迪奥斯圣母堂区，基谢拉莫宾河畔。\\
\noindent\textbf{保存状况:} 纸张受潮渍染；下缘有一种不明深色物质，穿透了整卷所有纸页。化学分析无定论。\\
\noindent\textbf{附入者:} 卢西亚·门德斯医生，2026年2月12日，附注：\textquotedblleft{\textit{这证明日记不是开端。日记只是回声。}\textquotedblright{}}

\noindent 省长大人：

\noindent 请允许我冒昧写信给您，因为活着的人已经不再听我说话，而死者太多，我一个人埋葬不完。

\noindent 自去年十一月起我便在此，受命为到达河岸的逃荒者施行圣事。一月份我接收了两百个灵魂。二月份两百四十个。这个三月，我已数不清了。他们成群结队地到来，整个家庭沦为行走的尘土。他们的眼睛是第一个死去的部分。先于身体。眼睛失去光、失去水、失去生命。它们睁着，却看不见。仿佛旱灾从眼睛里灌进去，从里面把灵魂晒干了。

\noindent 死者被运到距堂区半里格的一条沟坑。我没有祝圣过的墓地。没有做棺材的木头。没有裹尸布。我只有一条长长的土坑，开在干裂的土地上，尸体一层层地放进去。土盖上去。更多尸体。更多土。秃鹫学会了路。它们在天空围成一个永不散去的圆圈。夜里也不散。它们栖息在周围的树上，等待着。

\noindent 我要写信给您的正是这件事。

\noindent 三个夜晚以来，我无法安心为亡者诵经而不感到一种在场。不是灵性的在场——三十年的神职生涯让我习惯分辨那种在场。这是身体性的在场。它带着气味到来。

\noindent 气味来自沟坑。

\noindent 不是我所熟知的那种寻常死亡的气味。这是一种甜味。比乳香更甜，比仙人掌花更甜。一种从鼻孔进入、像浓稠液体一样滞留在喉咙里的气味。逃荒者也闻得到。他们转向沟坑的方向，一直望着。不哭。不祈祷。只是望着。仿佛在等什么东西从那里出来。

\noindent 第二天夜里，我亲自去了沟坑。

\noindent 我带了一盏提灯、圣水、一尊十字架。神父的东西。一直管用的东西。我走到坑边，照亮了坑底。

\noindent 我看见的不是尸体。

\noindent 是残骸。手臂和腿像根须一样从土里伸出来。张开的双手朝上，仿佛在乞求。一只孩子的手，小小的，手指扭曲，指着夜空，仿佛在指控某个人。

\noindent 那里的气味更浓。

\noindent 我跪下为沟坑祝圣。我开始念《永恒的安息》。念到第三个词，我感到土地震颤了。

\noindent 不是震动。是搏动。仿佛大地有一颗心脏。

\noindent 我跪着。把右手按在土上。土是温的。不是太阳晒的温热——是从里面透出的温热。病人身体的温度。土地在发烧。

\noindent 我把手插进土里。直到手腕。土是松软的、湿润的——在这干涸的荒野上是不可能的。我感到有什么东西擦过我的手指。光滑的。柔软的。某种在触碰时移动的东西，仿佛它是活的。

\noindent 我抽回了手。

\noindent 土自行合拢，封住了那个洞。

\noindent 附近一个男人——一个名叫维森特的逃荒者，上周失去了妻子和三个孩子——看着我说：

\noindent \textquotedblleft{\textit{神父先生，沟坑饿了。}\textquotedblright{}}

\noindent 大人，我知道您读到此处会作何想。认为孤绝与饥饿动摇了我的是非判断。认为酷热与死亡搅浑了我的理智。也许是这样。也许是。但我必须记录我所见，以便——如果我死在这里（我想我会死在这里）——有人知道这条沟坑不是寻常的沟坑。

\noindent 第三天——也就是昨天——我醒来时右手掌心有一道印记。

\noindent 一道深色圆形印记，在掌心中央。像烫伤，但不疼。皮肤完好。印记只是……在那里。仿佛有人在熟睡中把一枚滚烫的印章按在我的肉上。

\noindent 我试着洗。洗不掉。

\noindent 我试着祈祷。印记不消失。

\noindent 也许只是渗进皮肤的泥土。也许是正在开始的溃疡。也许是将要离去的理智。

\noindent 但今夜，当我闭上眼睛时，我听见了一个声音。

\noindent 不是来自外面的声音。是来自\textit{大地深处}的声音。它不停地重复同一个词，像正在学说话的人：

\noindent \textquotedblleft{\textit{饿。}\textquotedblright{}}

\noindent \textquotedblleft{\textit{饿。}\textquotedblright{}}

\noindent \textquotedblleft{\textit{饿。}\textquotedblright{}}

\noindent 一个孩子的声音。但没有一个孩子活下来开口说话。

\noindent 请允许我结束这封信。提灯的光在暗下去。

\noindent 愿上帝怜悯我们。

\noindent \textit{若泽·伊纳西奥·德苏萨神父}\\
\noindent \textit{雷梅迪奥斯圣母堂区}\\
\noindent \textit{1853年3月14日}

\noindent \textit{塞阿拉公共档案馆注：苏萨神父的信于1916年被人发现，夹在堂区的一本登记簿里，当时正值大楼翻修。神父于1853年3月15日去世——即写完这封信的第二天。死因登记为“自愿绝食”：神父停止进食，把自己的口粮让给了逃荒者。他的遗体被埋在他所描述的那条沟坑里。}

\noindent 本文件附有1916年的编目卡，上面一位匿名职员用铅笔在页边写道：\textquotedblleft{\textit{那股甜味从未离开过这里。转交米纳斯吉拉斯州公共档案馆。案卷1,260。}该职员身份不明。这段批注与若阿金日记页边出现的颤抖笔迹相同。\textquotedblright{}}

\noindent \textbf{2026年7月 —— 贝洛奥里藏特}\\
\noindent \textbf{卢西亚·门德斯医生最后的职业记录}

\noindent \textbf{内部文书 --- 自我转介申请临床离岗}\\
\noindent \textbf{专业人员:} 卢西亚·门德斯，CRM-MG 28.391。 \textbf{日期:} 2026年7月28日。

\noindent \textbf{伦理说明:}

\noindent 本记录不构成有效的心理鉴定。它是为证明立即离岗、终止研究并转介独立评估而写的。因属紧急记录，故保留第一人称。

\noindent \textbf{简要经过:}
\noindent 2026年1月，该专业人员在进行历史法医心理学研究时接触了归属于患者1,260若阿金·安东尼奥·科雷亚的日记。该材料间歇性保管，此后出现过封条与登记瑕疵。

\noindent 随后的数周，在接触档案后观察到持续存在的体征：

\begin{itemize}
\item 右眼颜色主观改变，尚无结论性眼科检查；
\item 反复感知到无环境来源的甜味；
\item 反复梦见沟坑、病房和空白表格；
\item 夜间过度警觉与注意力难以集中；
\item 即使在正式决定终止分析后仍有返回日记的冲动。
\end{itemize}

\noindent \textbf{工作假设:}

\begin{itemize}
\item 对研究创伤材料的急性反应，伴侵入性意象与失眠；
\item 长期文献接触强化的暗示作用；
\item 与患者1,260的反移情认同；
\item 轻度分离现象，有待外部评估确认。
\end{itemize}

\noindent \textbf{已采取的措施:}

\begin{itemize}
\item 暂停访谈、学术写作与副本流通；
\item 将原件转交司法保管；
\item 请求由未接触过本案的专业人员进行临床评估；
\item 建议不要单独阅读日记。
\end{itemize}

\noindent \textbf{未登记批注:}
\noindent 本文件的打印副本末尾，在专业人员姓名旁边有铅笔书写的数字“1.262”。笔迹与本人签名不符。其下方有一行空白，无身份标识。

\noindent \textit{卢西亚·门德斯 --- CRM-MG 28.391}\\
\noindent 记录于3点14分结束，按副本页脚所记时间。

\vspace{0.5em}
\noindent\textit{\textbf{↪ 路线:} \rota{DOC-01} \rotasep \rota{LUC-02} \rotasep \rota{EPI-01}}
